%% Separately authored historical derivative for the distinct Choice supplements.
%% Necessity: exact source lines 10559--10674 include Hartogs and the four-clause
%% Choice cycle, but do not contain the broader Zorn refinements or binary-tree
%% compactness proofs retained here.
%% Definition: this file is based on historical annex slice 17--79; the first
%% item is completed and retargeted to the bounded Choice-to-Zorn theorem.
%% Standing: this is a reciprocal editorial proof annex, not inherited source.
%% Nonredundancy: historical Hartogs lines 1--15 are retired as source overlap;
%% the retained slice proves only the separately routed additional results.
%% Use: annex/p06-route.tex inputs this file once, and its three annex identities
%% are reciprocal targets of the retained body results.

% BEGIN HISTORICAL SLICE 17-79
\begin{annexitem}{annex:p6.choice.1}{thm:p06-ac-to-zorn}{Choice implies Zorn's lemma}
\begin{proof}
Work in \(\mathsf{ZF}+\mathsf{AC}\).  Let \((P,\leq)\) satisfy the chain
hypothesis.  The empty chain has an upper bound, so \(P\ne\varnothing\).
Fix \(p_*\in P\).  By Theorem~\ref{thm:p01-well-ordering}, fix a well-order
\(\triangleleft\) of \(P\).  Write \(x<y\) for
\(x\leq y\land x\ne y\), and let \(h(P)\) be the Hartogs ordinal of \(P\).

Use transfinite recursion on the set ordinal \(h(P)\) to construct pairs
\((a_\alpha,x_\alpha)\in\{0,1\}\times P\).  Suppose the earlier pairs have
been constructed.  If every earlier flag is \(1\) and the set
\[
  C_\alpha=\{x_\beta:\beta<\alpha\}
\]
has a strict upper bound in \(P\), set \(a_\alpha=1\) and let
\(x_\alpha\) be the \(\triangleleft\)-least strict upper bound of
\(C_\alpha\).  Otherwise set \(a_\alpha=0\) and
\(x_\alpha=p_*\).  The defining operation is single-valued.  At each stage
Replacement collects the earlier pairs and their second coordinates, so
\(C_\alpha\) is a set; Separation forms the set of its strict upper bounds,
and the fixed well-order selects its least member.  Thus the set-recursion
theorem supplies the sequence on all of \(h(P)\).

If every flag were \(1\), then for \(\beta<\alpha<h(P)\) the element
\(x_\alpha\) would be strictly above \(x_\beta\).  Hence
\(\alpha\mapsto x_\alpha\) would be an injection \(h(P)\to P\),
contrary to Hartogs's theorem.  Let \(\alpha<h(P)\) be the least stage with
\(a_\alpha=0\).  Every earlier flag is \(1\), so \(C_\alpha\) is a chain,
and the definition of the stopping stage says that it has no strict upper
bound.  (The stage is nonzero because every member of nonempty \(P\) is a
strict upper bound of the empty chain.)

By hypothesis, \(C_\alpha\) has an upper bound \(u\in P\).  If
\(u\notin C_\alpha\), then every \(c\in C_\alpha\) satisfies
\(c<u\), making \(u\) a strict upper bound.  Therefore
\(u\in C_\alpha\) and it is the greatest member of that chain.  If some
\(y\in P\) satisfied \(u<y\), then transitivity would make \(y\) a strict
upper bound of all of \(C_\alpha\), again impossible.  Thus no element of
\(P\) lies strictly above \(u\), so \(u\) is maximal.
\end{proof}
\end{annexitem}

\begin{annexitem}{annex:p6.choice.2}{thm:p06-konig-infinity}{K\"onig's infinity lemma for the binary tree}
\begin{proof}
Call a finite binary sequence \(\mathbf s\) productive iff
\(\mathbf s\in\mathcal T\) and, for every
\(m\in\N\), it has an extension in \(\mathcal T\) of length
\(\operatorname{len}(\mathbf s)+m\).  The empty sequence is productive because
\(\mathcal T\) has a node at every level.  If a productive \(\mathbf s\) had neither
\(\mathbf s^\frown0\) nor \(\mathbf s^\frown1\) productive, each child would have a finite
bound on the heights of its extensions: this is immediate if the child is not
in \(\mathcal T\), and otherwise follows from the failure of productivity.  The maximum
of those two bounds would bound the height above \(\mathbf s\), a contradiction.
Recursively extend a productive node by \(0\) when its
\(0\)-child is productive, and by \(1\) otherwise.  This is a definition by
the fixed two-element order, not a sequence of arbitrary choices.  The union
of the recursively defined nodes is an \(f\in2^\omega\), and every
\(f\mathord{\upharpoonright}n\) lies in \(\mathcal T\); hence \(f\) is a branch.
\end{proof}
\end{annexitem}

\begin{annexitem}{annex:p6.choice.3}{thm:p06-classical-fan}{Classical fan theorem}
\begin{proof}
Suppose no uniform \(N\) exists.  Let \(\mathcal T\) consist of the binary sequences
none of whose initial segments belongs to \(B\).  This is a binary tree, and
the failure of a uniform bound supplies a node at every level.  By
Theorem~\ref{thm:p06-konig-infinity}, \(\mathcal T\) has a branch \(f\).  No initial
segment of \(f\) belongs to \(B\), contradicting that \(B\) is a bar.
\end{proof}
\end{annexitem}
% END HISTORICAL SLICE 17-79
